\documentclass[10pt, a4paper]{article}

% Parametri che modificano il file main.tex
% Le uniche parti da cambiare su main.tex sono:
% - vari \vspace tra sezioni
% - tabella azioni da intraprendere
% - sezione altro

\def\data{2023-11-17}
\def\oraInizio{09:00}
\def\oraFine{10:00}
\def\luogo{Piattaforma Discord}

\def\tipoVerb{Interno} % Interno - Esterno

\def\nomeResp{Feltrin E.} % Cognome N.
\def\nomeVer{Campese M.} % Cognome N.
\def\nomeSegr{Ciriolo I.} % Cognome N.

\def\nomeAzienda{Azzurro Digitale}
\def\firmaAzienda{azzurrodigitale.png}
\def\firmaResp{emanuele.png} % nome Responsabile

\def\listaPartInt{
Bresolin G.,
Campese M.,
Ciriolo I.,
Dugo A.,
Feltrin E.,
Michelon R.,
Orlandi G.
}

\def\listaPartEst{
Azzurro B.,
Digitale C.,
}

% Se nessuna revisione: \def\listaRevisioneAzioni {x}
\def\listaRevisioneAzioni {
{Feedback sull'assegnazione ruoli (avvenuta in precedenza) da parte di tutti i membri del team;},
{Controllo dei documenti in produzione con particolare attenzione verso l'impostazione del Piano di qualifica, Piano di progetto ed Analisi dei requisiti;},
{Verifica dei verbali prodotti durante la settimana.}
}

\def\listaOrdineGiorno {
{Assegnazione ruoli riunione;},
{Aggiornamento sul lavoro svolto da parte del team di amministratori (Bresolin G., Dugo A., Michelon R.) ed analisti (Ciriolo I., Orlandi G.);},
{Assegnazione task per la settimana successiva;},
{Breve spiegazione tenuta da Dugo A. relativa all'utilizzo del software \textit{Jira};},
{Presa in esame dell'idea proposta dall'azienda AzzurroDigitale di utilizzare di un \textit{repository} messo a disposizione da loro.}
}

\def\listaDiscussioneInterna {
{Assegnazione ruoli riunione:
    \begin{itemize}
        \item Responsabile: Feltrin E.;
        \item Segretario di riunione: Ciriolo I.;
        \item Responsabile della revisione: Campese M.
    \end{itemize}},
{Durante l'aggiornamento, il team di amministratori presenta alcuni cambiamenti effettuati nella gestione e versionamento della documentazione; segue la presentazione delle bozze dei diagrammi dei Casi d'uso da parte di Ciriolo I.;},
{A seguito di una breve discussione sono state assegnate le azioni da intraprendere che si trovano nell'apposita sezione a pagina seguente;},
{Successivamente alla spiegazione sull'utilizzo di \textit{Jira} si è pensato di utilizzarlo per l'organizzazione del lavoro del team. Questo software permetterebbe di specificare le dipendenze nei \textit{diagrammi di Gantt} e imporrebbe di chiudere tutti i task di una \textit{milestone} prima che questa possa essere considerata raggiunta;},
{Si è discussa la possibilità di utilizzare il \textit{repository} offerto dall'azienda proponente, notando però che non verrebbero stanziate particolari agevolazioni per il team, che dispone già di un \textit{repository} avviato.}
}

\newcommand{\domris}[2]{\textbf{#1}\\#2}

\def\listaDiscussioneEsterna {
\domris
{Domanda 1
}
{Riposta 1;
},
\domris
{Domanda 2
}
{Risposta 2;
},
\domris
{Domanda 3
}
{Risposta 3.
}
}

% Se nessuna decisione: \def\listaDecisioni {x}
\def\listaDecisioni {
{Si è deciso definitivamente di non utilizzare il \textit{repository} messo a disposizione dall'azienda proponente;},
{A seguito di una consultazione tra gli amministratori, il team utilizzerà il software \textit{Jira} a partire da questa riunione in avanti come \textit{ITS} ufficiale del gruppo, il quale sostituirà quello precedentemente utilizzato su \textit{GitHub}. Lo spazio di lavoro sarà organizzato dagli amministratori Dugo A. e Bresolin G.;},
{Soppressione \textit{issue} per la creazione di template per le \textit{issue} stesse in seguito al passaggio su \textit{Jira};},
{I diagrammi dei Casi d'uso saranno prodotti da Ciriolo I. e revisionati da Bresolin G.;},
{Il nome del repository utilizzato dal team sarà \textit{Knowledge Managment AI}, nome del capitolato dell'azienda proponente;},
{Durante il weekend avrà luogo un workshop presentato dal team di amministratori sulla configurazione dell'ambiente di lavoro e sul funzionamento degli \textit{script} in \textit{Python}.}
}

\usepackage{style}
\usepackage{headerfooter}
\usepackage{caption}
\usepackage{hyperref}
\usepackage{svg}
\usepackage{comment}
\usepackage{setspace}
\definecolor{primarycolor}{RGB}{248, 182, 143}

\title{\titolo}
\author{SWEetCode}

\begin{document}

% PRIMA PAGINA
\begin{titlepage}
    \thispagestyle{empty}
    \begin{tikzpicture}[remember picture, overlay]
        % TRIANGOLI
        \draw[fill=secondarycolor, secondarycolor] (current page.north west) -- (current page.south west) -- (8.8, -28);
        \draw[fill=primarycolor, primarycolor] (-3, 5) -- (4, -13.6) -- (11, 5);

        % LOGO
        \node [xshift=-5cm, yshift=25cm] (logo) at (current page.south east) {\includegraphics[width=6.5cm]{img/logo.png}};

        % SWEETCODE - DATE
        \node [anchor=north east, align=right, xshift=-1.2cm, yshift=20.5cm, text=black] (sweetcode) at (current page.south east) {\fontsize{32pt}{36pt}\selectfont SWEetCode};
        \draw[line width=4pt, lightcol] ([xshift=-3cm, yshift=-0.37cm]sweetcode.south west) -- ([yshift=-0.37cm]sweetcode.south east);
        \node [anchor=north east, align=right, xshift=-1.2cm, yshift=18.7cm, text=black] (date) at (current page.south east){\fontsize{24pt}{24pt} \selectfont Verbale \tipoVerb};

        % NOME FILE
        \node [anchor=north east, text width=15cm, align=right, xshift=-1.2cm, yshift=17cm, text=black] (titolo) at (current page.south east){\fontsize{48pt}{48pt}\textbf{\data}};

        % BOX DATI PARTECIPANTI
        \ifthenelse{\equal{\tipoVerb}{Esterno}}{
            \node[anchor=north east, xshift=-1.2cm, yshift=14.5cm, minimum width=8cm] (box) at (current page.south east){};
        }{
            \node[anchor=north east, xshift=-1.2cm, yshift=12.5cm, minimum width=8cm] (box) at (current page.south east){};
        }
 
        % RESPONSABILE
        \node[anchor=north west, align=left] (dati2) at (box.north west) {\fontsize{15pt}{15pt}\selectfont \textbf{Responsabile}};
        \draw[line width=4pt, lightcol] (dati2.south west) -- ([xshift=8cm]dati2.south west);
        \node[anchor=north west, align=left] (dati21) at (dati2.south west){\fontsize{13pt}{13pt}\selectfont \nomeResp};

        % VERIFICATORE
        \node[anchor=north west, yshift=-1cm, align=left] (dati3) at (dati21.north west) {\fontsize{15pt}{15pt}\selectfont \textbf{Verificatore}};
        \draw[line width=4pt, lightcol] (dati3.south west) -- ([xshift=8cm]dati3.south west);
        \node[anchor=north west, align=left] (dati31) at (dati3.south west){\fontsize{13pt}{13pt}\selectfont \nomeVer};

        % SEGRETARIO DI RIUNIONE
        \node[anchor=north west, yshift=-1cm, align=left] (dati4) at (dati31.north west) {\fontsize{15pt}{15pt}\selectfont \textbf{Segretario di Riunione}};
        \draw[line width=4pt, lightcol] (dati4.south west) -- ([xshift=8cm]dati4.south west);
        \node[anchor=north west, align=left] (dati41) at (dati4.south west){\fontsize{13pt}{13pt}\selectfont \nomeSegr};
        
        % UNIPD - SWE
        \node [xshift=4.4cm, yshift=2.3cm, draw, secondarycolor, text=white] (uni) at (current page.south west) {\fontsize{20pt}{20pt} \selectfont Università di Padova};
        \node [xshift=0.65cm, yshift=0.7cm, draw, secondarycolor, text=white, below=of uni] (corso) {\fontsize{20pt}{20pt}\selectfont Ingegneria del Software};

        % FIRMA AZIENDA
        \ifthenelse{\equal{\tipoVerb}{Esterno}}{
            \draw[line width=4pt, lightcol] ([xshift=-1.2cm, yshift=5.2cm]current page.south east) -- ([xshift=-8cm, yshift=5.2cm]current page.south east);
            \node [xshift=-4.8cm, yshift=6cm] (logo) at (current page.south east) {\includegraphics[width=6cm]{img/firme/\firmaAzienda}};
            \node[anchor=north west, xshift=12.9cm, yshift=4.95cm, align=left] at (current page.south west)
            {\fontsize{13pt}{13pt}\selectfont L'azienda: \nomeAzienda};
        }     

        % FIRMA
        \draw[line width=4pt, lightcol] ([xshift=-1.2cm, yshift=1.8cm]current page.south east) -- ([xshift=-8cm, yshift=1.8cm]current page.south east);
        \node [xshift=-4.8cm, yshift=2.45cm] (logo) at (current page.south east) {\includegraphics[width=6cm]{img/firme/\firmaResp}};
        \node[anchor=north west, xshift=12.9cm, yshift=1.45cm, align=left] at (current page.south west)
        {\fontsize{13pt}{13pt}\selectfont Il Responsabile: \nomeResp};
        
    \end{tikzpicture}
\end{titlepage}

% REGISTRO DELLE VERSIONI
%Registro in ordine dalla più recente alla meno recente!
{\renewcommand{\arraystretch}{1.5}
\section*{Registro delle versioni}

\begin{xltabular}{\textwidth}{c|c|c|c|X}
\label{tab:long}

\textbf{Versione} & \textbf{Data} & \quantities{\textbf{Responsabile di}\\\textbf{stesura}}& \textbf{Revisore} & \quantities{\textbf{Dettaglio e}\\\textbf{motivazioni}} \\
\endfirsthead

\textbf{Versione} & \textbf{Data} & \quantities{\textbf{Responsabile di}\\\textbf{stesura}}& \textbf{Revisore} & \quantities{\textbf{Dettaglio e}\\\textbf{motivazioni}} \\
\endhead

\multicolumn{5}{r}{{Continua nella pagina successiva}} \\
\endfoot

\endlastfoot

\hline
v3.0.0(0) & $2024-04-12$ & - & - & Approvazione versione finale PB.\\
\hline
v2.27.9(3) & $2024-04-12$ & \quantities{Dugo A.} & Michelon R. & Consuntivo Sprint11 e finale.\\
\hline
v2.27.7(15) & $2024-04-06$ & \quantities{Bresolin G.} & Dugo A. & Correzione consuntivo totale Sprint 10.\\
\hline
v2.27.7(4) & $2024-03-30$ & \quantities{Bresolin G.} & Feltrin E. & Consuntivo Sprint10 e aggiornamento pianificazione Sprint11.\\
\hline
v2.27.1(1) & $2024-03-26$ & \quantities{Orlandi G.} & Bresolin G. & Consuntivo Sprint9 e aggiornamento pianficazione Sprint10 e Sprint11.\\
\hline
v2.0.0(37) & $2024-03-02$ & \quantities{Ciriolo I.} & Bresolin G. & Consuntivo Sprint8, aggiornamento pianificazione Sprint9 e Sprint10, aggiornamento analisi dei rischi e mitigazione.\\
\hline
v2.0.0(2) & $2024-02-20$ & \quantities{Campese M.} & Ciriolo I. & Correzione della organizzazione delle sezioni.\\
\hline
v2.0.0(0) & $2024-02-11$ & - & - & Approvazione versione finale RTB.\\
\hline
v1.10.4(18) & $2024-02-11$ & \quantities{Campese M.} & Michelon R. & Aggiornamenti di riferimenti e modifiche correttive.\\
\hline
v1.10.4(16) & $2024-02-11$ & \quantities{Bresolin G.} & Michelon R. & Consuntivo Sprint7 e pianificazione Sprint8 e Sprint9.\\
\hline
v1.10.4(15) & $2024-02-10$ & \quantities{Orlandi G.} & Michelon R. & Consuntivo Sprint6 e pianificazione Sprint7 e Sprint8.\\
\hline
v1.10.4(11) & $2024-02-08$ & \quantities{Michelon R.} & Bresolin G. & Consuntivo Sprint5 e aggiornamento pianificazione Sprint6 e Sprint7.\\
\hline
v1.10.1(21) & $2024-01-08$ & \quantities{Dugo A.} & Michelon R. & Refactoring preventivi, consuntivo Sprint4 e aggiornamento pianificazione Sprint5 e Sprint6.\\
\hline
v1.10.1(18) & $2024-01-05$ & \quantities{Dugo A.} & Michelon R. & Consuntivo Sprint3 e aggiornamento pianificazione Sprint4 e Sprint5.\\
\hline
v1.10.1(5) & $2024-01-02$ & \quantities{Dugo A.} & Michelon R. & Consuntivo Sprint2 e aggiornamento pianificazione Sprint3 e Sprint4.\\
\hline
v1.10.1(3) & $2023-12-31$ & \quantities{Dugo A.} & Bresolin G. & Consuntivo Sprint1 e aggiornamento pianificazione Sprint2 e Sprint3.\\
\hline
v1.0.0(9) & $2023-11-27$ & \quantities{Campese M.} & Ciriolo I. & Prima stesura, sezione rischi e sezione primo periodo.
\end{xltabular}
}
%\newpage
\newpage
%todo modificare questa section per avere sempre un titolo "grande" ma senza numero
\section*{Candidatura per Product Baseline}
Il gruppo SWEetCode è lieto di presentare la propria candidatura per la revisione di avanzamento Product Baseline(PB), nell'impegno di consegnare il progetto da Voi commissionato, denominato \textbf{Knowledge Management AI} (C1), proposto dall'azienda \textbf{AzzurroDigitale}.\\

Per una panoramica delle attività e una consultazione semplificata dei documenti riguardanti l'PB, il team di SWEetCode estende un invito cordiale a visitare il sito:
\begin{center}
\href{https://sweetcode-team.github.io}{\textcolor{blue}{sweetcode-team.github.io}}
\end{center}

\vspace{1em}
In tale spazio si potranno trovare:
\begin{itemize}
    \item (Analisi dei requisiti v3.0.0(0));
    \item (Glossario v3.0.0(0));
    \item (Norme di progetto v3.0.0(0));
    \item (Piano di progetto v3.0.0(0));
    \item (Piano di qualifica v3.0.0(0));
    \item (Manuale utente v3.0.0(0));
    \item I verbali delle riunioni interne tenute i giorni:
        \begin{itemize}
            \item 2024-02-19;
            \item 2024-02-26;
            \item 2024-03-07;
            \item 2024-03-12;
            \item 2024-03-19;
            \item 2024-03-25;
            \item 2024-04-01;
            \item 2024-04-12.
        \end{itemize}
   \item I verbali delle riunioni esterne tenute i giorni:
        \begin{itemize}
            \item 2023-02-28: AzzurroDigitale;
            \item 2023-03-13: AzzurroDigitale;
            \item 2023-03-26: AzzurroDigitale.
        \end{itemize}    
\end{itemize}

Per consultare il repository del progetto, il team SWEetCode estende un invito cordiale a visitare il link:\\
\begin{center}
\href{https://github.com/sweetcode-team/Knowledge_Management_AI}{\textcolor{blue}{github.com/sweetcode-team/Knowledge\_Management\_AI}}
\end{center}

\newpage
\subsection*{Aggiornamento Impegni}
Il team conclude con la Product Baseline lo sviluppo del progetto confermando la data di consegna prevista per il 2024-04-12.
Il team ha rispettato il preventivo iniziale di 12715,00€, arrivando ad un consumo totale di 12450,00€.
Le tabelle orarie sono riportate nel documento (Piano di progetto v3.0.0(0), \S Stima dei costi di realizzazione).

{\renewcommand{\arraystretch}{1.5}
\begin{table}[H]
\centering
\begin{tabularx}{0.54\textwidth}{c|c|c}
\rowcolor{primarycolor}
    \textbf{Revisione} & \textbf{Data}& \textbf{Costi}\\
\hline
    \textbf{Product Baseline} & 2024-04-12 & 12450,00€\\
\end{tabularx}
\end{table}

\vspace{1em}
Il team è composto dai seguenti membri:
\\
\\
% TABELLA (tenere anche arraystretch che aumenta il padding verticale)
{\renewcommand{\arraystretch}{1.5}
\begin{tabularx}{\textwidth}{X|X}
\textbf{Cognome e Nome} & \textbf{Matricola n°} \\
\hline
Bresolin Gianluca & 2034316 \\
\hline
Campese Martina & 1122439 \\
\hline
Ciriolo Irene & 2043682 \\
\hline
Dugo Alberto & 2042382 \\
\hline
Feltrin Emanuele & 2034314 \\
\hline
Michelon Riccardo & 2042341 \\
\hline
Orlandi Giacomo & 1122253 \\
\end{tabularx}}\\

Cordiali saluti,\\

Il gruppo SWEetCode

\end{document}